% GENERATED FILE. Edit canonical lessons, then run npm run generate:books.
% canonical-source-hash: fnv1a64:3b24e8ff3fa90908
% canonical-lessons: SA-C292-carah, SA-C292-pramanam, SA-C292-pariksate, SA-C292-rahasyam, SA-C292-prakatayati

\chapter{The Investigation}
\label{ch:sa-the-investigation}
\hlchaptermodality{292}

\begin{chapteropening}
\textbf{By the end of this chapter:} \emph{I can say what the investigator found and what happens if there is proof.}

Complete the last lesson of chapter 292: I can say what the investigator found and what happens if there is proof.
\end{chapteropening}

\section[\texorpdfstring{\sk{चरः} (caraḥ)}{चरः (caraḥ)}]{\sk{चरः} (caraḥ) — a spy; a secret agent}

\label{lesson:SA-C292-carah}



\begin{quote}
Before the new one: say the Sanskrit for fraud, then the Sanskrit for fake.
\end{quote}

\subsection*{You'll want to know: \sk{चरः}}
\textbf{\sk{चरः}} — \emph{caraḥ} — \textquotedblleft{}a spy; a secret agent\textquotedblright{}.

\begin{quote}
\sk{राजा} \sk{चरं} \sk{प्रेषयति।}
\end{quote}

\begin{quote}
\emph{rājā caraṃ preṣayati.}
\end{quote}

\begin{quote}
The king sends a spy.
\end{quote}

A masculine noun.

\begin{cousinweb}
\textbf{\sk{चरः}} (\emph{caraḥ}) comes from the root √car, \textquotedblleft{}to move, go about\textquotedblright{}; MW gives \emph{cara} as a spy or secret agent (MW).
\end{cousinweb}

\begin{grammarlens}[title={What you've built}]
One more word for telling what happened.
\end{grammarlens}

\subsection*{Your turn}
\begin{itemize}
\raggedright
  \item \emph{Say it:} \emph{caraḥ}
  \item \emph{Say it:} \emph{caraḥ}, once more
  \item \emph{Say it:} \emph{rājā caraṃ preṣayati.} — The king sends a spy.
  \item \emph{Recall:} say the Sanskrit for fraud, then the Sanskrit for fake, then say \emph{caraḥ} again
\end{itemize}

\begin{tcolorbox}[breakable,colback=teal!4,colframe=teal!35!black,title={Before you move on}]
What does \sk{चरः} mean? (\textquotedblleft{}A spy; a secret agent\textquotedblright{}.) Say it once more.
\end{tcolorbox}
\section[\texorpdfstring{\sk{प्रमाणम्} (pramāṇam)}{प्रमाणम् (pramāṇam)}]{\sk{प्रमाणम्} (pramāṇam) — proof; evidence}

\label{lesson:SA-C292-pramanam}



\begin{quote}
Before the new one: say the Sanskrit for fake, then the Sanskrit for a spy.
\end{quote}

\subsection*{You'll want to know: \sk{प्रमाणम्}}
\textbf{\sk{प्रमाणम्}} — \emph{pramāṇam} — \textquotedblleft{}proof; evidence\textquotedblright{}.

\begin{quote}
\sk{चरः} \sk{प्रमाणम्} \sk{अन्विष्यति।}
\end{quote}

\begin{quote}
\emph{caraḥ pramāṇam anviṣyati.}
\end{quote}

\begin{quote}
The spy searches for proof.
\end{quote}

A neuter noun.

\begin{cousinweb}
\textbf{\sk{प्रमाणम्}} (\emph{pramāṇam}) is built from \emph{pra} and the root √mā, \textquotedblleft{}to measure\textquotedblright{}: a \emph{pramāṇa} is first a measure or standard, and MW also gives it as any proof or evidence (MW).
\end{cousinweb}

\begin{grammarlens}[title={What you've built}]
One more word for telling what happened.
\end{grammarlens}

\subsection*{Your turn}
\begin{itemize}
\raggedright
  \item \emph{Say it:} \emph{pramāṇam}
  \item \emph{Say it:} \emph{pramāṇam}, once more
  \item \emph{Say it:} \emph{caraḥ pramāṇam anviṣyati.} — The spy searches for proof.
  \item \emph{Recall:} say the Sanskrit for fake, then the Sanskrit for a spy, then say \emph{pramāṇam} again
\end{itemize}

\begin{tcolorbox}[breakable,colback=teal!4,colframe=teal!35!black,title={Before you move on}]
What does \sk{प्रमाणम्} mean? (\textquotedblleft{}Proof; evidence\textquotedblright{}.) Say it once more.
\end{tcolorbox}
\section[\texorpdfstring{\sk{परीक्षते} (parīkṣate)}{परीक्षते (parīkṣate)}]{\sk{परीक्षते} (parīkṣate) — examines}

\label{lesson:SA-C292-pariksate}



\begin{quote}
Before the new one: say the Sanskrit for a spy, then the Sanskrit for proof.
\end{quote}

\subsection*{You'll want to know: \sk{परीक्षते}}
\textbf{\sk{परीक्षते}} — \emph{parīkṣate} — \textquotedblleft{}examines\textquotedblright{}.

\begin{quote}
\sk{चरः} \sk{चिह्नं} \sk{परीक्षते।}
\end{quote}

\begin{quote}
\emph{caraḥ cihnaṃ parīkṣate.}
\end{quote}

\begin{quote}
The spy examines the mark.
\end{quote}

Past (he …): \sk{सः} \sk{चिह्नं} \sk{परीक्षितवान्}, \emph{saḥ cihnaṃ parīkṣitavān}.

\begin{cousinweb}
\textbf{\sk{परीक्षते}} (\emph{parīkṣate}) is built from \emph{pari}, \textquotedblleft{}around\textquotedblright{}, and the root √īkṣ, \textquotedblleft{}to see\textquotedblright{}; MW gives \emph{parīkṣ} as to look round, inspect carefully and examine (MW).
\end{cousinweb}

\begin{grammarlens}[title={What you've built}]
One more word for telling what happened.
\end{grammarlens}

\subsection*{Your turn}
\begin{itemize}
\raggedright
  \item \emph{Say it:} \emph{parīkṣate}
  \item \emph{Say it:} \emph{parīkṣate}, once more
  \item \emph{Say it:} \emph{caraḥ cihnaṃ parīkṣate.} — The spy examines the mark.
  \item \emph{Recall:} say the Sanskrit for a spy, then the Sanskrit for proof, then say \emph{parīkṣate} again
\end{itemize}

\begin{tcolorbox}[breakable,colback=teal!4,colframe=teal!35!black,title={Before you move on}]
What does \sk{परीक्षते} mean? (\textquotedblleft{}Examines\textquotedblright{}.) Say it once more.
\end{tcolorbox}
\section[\texorpdfstring{\sk{रहस्यम्} (rahasyam)}{रहस्यम् (rahasyam)}]{\sk{रहस्यम्} (rahasyam) — a secret}

\label{lesson:SA-C292-rahasyam}



\begin{quote}
Before the new one: say the Sanskrit for proof, then the Sanskrit for examines.
\end{quote}

\subsection*{You'll want to know: \sk{रहस्यम्}}
\textbf{\sk{रहस्यम्}} — \emph{rahasyam} — \textquotedblleft{}a secret\textquotedblright{}.

\begin{quote}
\sk{चरः} \sk{रहस्यं} \sk{जानाति।}
\end{quote}

\begin{quote}
\emph{caraḥ rahasyaṃ jānāti.}
\end{quote}

\begin{quote}
The spy knows the secret.
\end{quote}

A neuter noun.

\begin{cousinweb}
\textbf{\sk{रहस्यम्}} (\emph{rahasyam}) is built on \emph{rahas}, \textquotedblleft{}a lonely place, secrecy\textquotedblright{}; MW gives \emph{rahasya} as secret and, as a noun, a secret or mystery (MW).
\end{cousinweb}

\begin{grammarlens}[title={What you've built}]
One more word for telling what happened.
\end{grammarlens}

\subsection*{Your turn}
\begin{itemize}
\raggedright
  \item \emph{Say it:} \emph{rahasyam}
  \item \emph{Say it:} \emph{rahasyam}, once more
  \item \emph{Say it:} \emph{caraḥ rahasyaṃ jānāti.} — The spy knows the secret.
  \item \emph{Recall:} say the Sanskrit for proof, then the Sanskrit for examines, then say \emph{rahasyam} again
\end{itemize}

\begin{tcolorbox}[breakable,colback=teal!4,colframe=teal!35!black,title={Before you move on}]
What does \sk{रहस्यम्} mean? (\textquotedblleft{}A secret\textquotedblright{}.) Say it once more.
\end{tcolorbox}
\section[\texorpdfstring{\sk{प्रकटयति} (prakaṭayati)}{प्रकटयति (prakaṭayati)}]{\sk{प्रकटयति} (prakaṭayati) — reveals}

\label{lesson:SA-C292-prakatayati}



\begin{quote}
Before the new one: say the Sanskrit for examines, then the Sanskrit for a secret.
\end{quote}

\subsection*{You'll want to know: \sk{प्रकटयति}}
\textbf{\sk{प्रकटयति}} — \emph{prakaṭayati} — \textquotedblleft{}reveals\textquotedblright{}.

\begin{quote}
\sk{यदि} \sk{प्रमाणं} \sk{भविष्यति}, \sk{तर्हि} \sk{सः} \sk{प्रकटयिष्यति।}
\end{quote}

\begin{quote}
\emph{yadi pramāṇaṃ bhaviṣyati, tarhi saḥ prakaṭayiṣyati.}
\end{quote}

\begin{quote}
If there is proof, then he will reveal it.
\end{quote}

Past (he …): \sk{सः} \sk{रहस्यं} \sk{प्रकटितवान्}, \emph{saḥ rahasyaṃ prakaṭitavān}.

Future: \sk{प्रकटयिष्यति}, \emph{prakaṭayiṣyati}.

\begin{cousinweb}
\textbf{\sk{प्रकटयति}} (\emph{prakaṭayati}) is built on \emph{prakaṭa}, \textquotedblleft{}evident, manifest\textquotedblright{}; MW gives \emph{prakaṭayati} as a verb made from it, to manifest or disclose (MW).
\end{cousinweb}

\begin{grammarlens}[title={What you've built}]
One more word for telling what happened.
\end{grammarlens}

\subsection*{Your turn}
\begin{itemize}
\raggedright
  \item \emph{Say it:} \emph{prakaṭayati}
  \item \emph{Say it:} \emph{prakaṭayati}, once more
  \item \emph{Say it:} \emph{yadi pramāṇaṃ bhaviṣyati, tarhi saḥ prakaṭayiṣyati.} — If there is proof, then he will reveal it.
  \item \emph{Recall:} say the Sanskrit for examines, then the Sanskrit for a secret, then say \emph{prakaṭayati} again
\end{itemize}

\begin{tcolorbox}[breakable,colback=teal!4,colframe=teal!35!black,title={Before you move on}]
What does \sk{प्रकटयति} mean? (\textquotedblleft{}Reveals\textquotedblright{}.) Say it once more.
\end{tcolorbox}
